Persistent homology reliably detects temporal recurrence structure in trajectories of neural audio embeddings. However, this structure corresponds only weakly to musical genre semantics, depends strongly on the choice of embedding model, and is localized on repeated musical sections only for a minority of tracks. On the GTZAN dataset (999 tracks, 10 genres), a reproducible pipeline is built over four representations (MERT, MuQ, EnCodec, MIR) with a ladder of three surrogate controls. Topological features classify genre near chance level (macro-F1~$\approx 0.22$), whereas time-averaged embeddings reach~$0.82$, exposing a sharp separation between trajectory topology and genre-level semantic organization. The work also identifies several methodological pitfalls, including an almost twofold overestimation of the share of ``meaningful'' loops caused by treating cocycle support as a geometric loop representative. The contribution is a map of where TDA is applicable to audio embeddings and a rigorous methodological template with controls, effect sizes, and reproducibility mechanisms; the negative results are a central contribution rather than a failure.
